% Section VI: Discussion and Conclusion. Numbers repeated from Section V
% (results/phase2/metrics/, results/phase2/tuning/ and results/metrics/).
\section{Discussion and Conclusion}\label{sec:conclusion}

\emph{What worked.} One shared network learns all three tasks with a third
of the weights of three separate networks. Its regularization matters:
without dropout and weight decay the CV error rises by 1.08~m, and the heard
indicators help position too (0.58~m); the effect of hidden WAPs (0.16~m) is at
the level of chance. The result hardly depends on the seed (test:
$8.21\pm0.06$~m over five seeds). The second hybrid, the blend with the
Random Forest, beat both models in CV (9.14~m) and improved the network on
the test set (8.10~m, 92.2\% of floors).

\emph{What did not.} Three findings go against our expectations, and we report
them as they are. First, the Phase~1 features did not help the network (a
change of 0.04~m, within chance); a likely reason, which we did not test, is
that with 18,719 training scans it learns from the readings alone where the
WAPs are. Second, sharing one trunk did not help accuracy: three single-task
networks were better in CV and on the test set (8.03 against 8.32~m, 94.1\%
against 90.8\% of floors). For position, a single-task network also beat a
SIMO-style shared network on this dataset, although a multi-task grid-cell
model beat both~\cite{laska2022multitask}. Third, user-group CV ranked the
network ahead of the Random Forest, but on the test set the forest was better
(7.42 against 8.32~m): CV with 16 user groups is a strict but noisy guide, and
its ranking of very different models need not hold on a test set recorded
months later.

\emph{Conclusion.} Phase~2 built the multi-task network of the brief on the
same data, folds, test set and metrics as Phase~1, with an ablation study and
seed checks. The hybrid MLP reached 99.9\% building accuracy, 90.8\% floor
accuracy and a mean 2D error of 8.32~m on the test set; its blend with the
Phase~1 Random Forest reached 8.10~m and 92.2\%, while the forest alone remains
the best on position (7.42~m). Learned
task weights~\cite{kendall2018multi} and inputs that are more robust to new
phones are the next steps.
